%%=====================================
%% A CV template made by Garland Culbreth (culbreth.garland@gmail.com)
%% (v1.0.0, 2021-09-04)
%%
%% It is a heavily modified derivative of altacv.cls v1.6.2
%% https://github.com/liantze/AltaCV
%%
%% It may be distributed and/or modified under the
%% conditions of the LaTeX Project Public License, either version 1.3
%% of this license or (at your option) any later version.
%% The latest version of this license is in
%%    http://www.latex-project.org/lppl.txt
%% and version 1.3 or later is part of all distributions of LaTeX
%% version 2003/12/01 or later.
%%
%%=====================================
%% !!!---Edit content in the .tex files in the "content" folder---!!!

\documentclass[11pt,letterpaper,ragged2e,withhyper]{dodiresume}


%% change gap between columns if you want
\geometry{left=0.5in,right=0.5in,top=0.2in,bottom=0.7in,columnsep=0.4in}

%% typesets columns of text in parallel
\usepackage{paracol}

% Used to customize footer.
\usepackage{fancyhdr}


%% uncomment and edit to change name and section color
% \definecolor{custom}{HTML}{582766}  % hex code
% \colorlet{name}{custom}
% \colorlet{heading}{custom}


%% uncomment to use small caps instead of all caps for sections
%% only works if the font supports small caps
% \renewcommand{\cvsectionfont}{\Large\scshape\bfseries}


%% uncomment and edit to change name font size
% \renewcommand{\namefont}{\bfseries\fontsize{36pt}{36pt}\selectfont}



\usepackage{helvet}  % Helvetica is always reliable
\renewcommand{\familydefault}{\sfdefault}

% Set footer.
\pagestyle{fancy}
\renewcommand{\headrulewidth}{0pt}
\fancyhf{}
\fancyfoot[L]{\small\textcolor{body}{Autorizzo il trattamento dei dati personali contenuti nel mio curriculum vitae in base all’art. 13 GDPR (Regolamento UE 2016/679) ai fini della ricerca e selezione del personale.}}

\begin{document}

% Document's header.
\name{Emanuele Petriglia}
\personalinfo{%
    %% Not all of these are required. More than 4 will look bad.
    \location{Aggiornato al 2023-02-09} \\
    \location{Domicilio a Milano} Età 24\\
    %\phone{+1 (123) 456-7890} \\
    \email{work@emanuelepetriglia.com} \\
    \linkedin{linkedin.com/in/emanuele-petriglia} \\
    %\homepage{www.emanuelepetriglia.com} \\
    \github{github.com/ema-pe} \github{gitlab.com/ema-pe}
    % \orcid{0000-0000-0000-0000}
}
\makecvheader  

%% makes itemize fonts size slightly smaller
\AtBeginEnvironment{itemize}{\small}  

% sets the left:right column width ratio to 65:35.
\columnratio{0.60}



%% Start a 2-column paracol. Both columns automatically break 
%% across pages if things get too long.
\begin{paracol}{2}

Studente magistrale di Informatica e attivista nel mondo del software libero e open source.

% -------------------
% Experience section.
% -------------------
\cvsection{Esperienza}

\cvevent{Sviluppatore Go, Progettista API e database relazionale}{Università degli Studi di Roma ``La Sapienza''}{feb 2021 -- set 2021}
\begin{itemize}\setlength\itemsep{0.2em}
\item Analisi di un servizio di messagistica domain-specific per SeismoCloud.
\item Progettazione di \textbf{API HTTP REST} con OpenAPI.
\item Progettazione e implementazione in \textbf{SQL} di un modello per il database relazionale \textbf{MariaDB}.
\item Implementazione delle API in linguaggio \textbf{Go} in un ambiente di sviluppo locale centrato su \textbf{Git} e GitLab supportato da Docker.
\item Il lavoro si è svolto in maniera indipendente con riunioni settimanali con il team per discussioni e revisione del lavoro.
\end{itemize}

\cvevent{Sviluppatore C, Sistemista UNIX}{INFN -- Laboratori Nazionali di Frascati, Roma}{giu 2016 -- lug 2016}
\begin{itemize}\setlength\itemsep{0.2em}
\item Utilizzo del linguaggio \textbf{C} per risolvere problemi, alcuni dei quali legati alla gestione dei dati provenienti dall'acceleratore di particelle KLOE.
\item Esplorazione e studio del sistema operativo \textbf{IBM AIX} e delle principali soluzioni e tecnologie adottate nel centro di calcolo.
\end{itemize}

% ------------------
% Education section.
% ------------------
\cvsection{Formazione}

\cvevent{Laurea Magistrale in Informatica}{Università degli Studi di Milano-Bicocca}{2021 -- in corso}

\cvevent{Laurea Triennale in Informatica}{Università degli Studi di Roma ``La Sapienza''}{2017 -- 2021}
\begin{itemize}\setlength\itemsep{0.2em}
\item Voto: 110/110 e lode
\end{itemize}

\cvevent{Diploma in Informatica}{ITIS ``G. Vallauri'', Velletri, Roma}{2012 -- 2017}
\begin{itemize}\setlength\itemsep{0.2em}
\item Voto: 94/100
\end{itemize}

% -----------------
% Projects section.
% -----------------
\cvsection{Progetti}

\cvproject{SapienzaHub}{mag 2019 -- apr 2022}{\github{https://gitlab.com/sapienzastudents}}{}
\begin{itemize}\setlength\itemsep{0.2em}
\item Progetto indipendente creato e gestito da un gruppo di studenti per riunire le informazioni legate all'università.
\item Progettazione e sviluppo di un sito Web statico e di un bot sviluppato in Go per la moderazione e gestione della rete di gruppi Telegram.
\end{itemize}

% ------------------
% Interests section.
% ------------------
\cvsection{Interessi}

\begin{itemize}\setlength\itemsep{0.2em}
\normalsize
\item Professionali: software libero e open source (FOSS), sviluppo di software backend e di sistema, protezione dei dati personali e GDPR.
\item Personali: sport (di gruppo e singoli), fotografia, scrittura di racconti, lettura di narrativa.
\end{itemize}

% ------------------
% Languages section.
% ------------------
\cvsection{Lingue}

\begin{itemize}\setlength\itemsep{0.2em}
\normalsize
\item Inglese (B1)
\end{itemize}

\switchcolumn

% ---------------
% Skills section.
% ---------------
\cvsection{Skills}
\vspace{-6pt}

\cvsubsection{Linguaggi}
\cvskill{Go, C, Python, Java, PHP, Scala, C++, SQL}

\vspace{-2mm}
\cvsubsection{Tecnologie}
\cvskill{Linux (Fedora, CentOS, Debian, Ubuntu), Git, GNU tools (Bash, m4, AWK, grep...)}

\vspace{-2mm}
\cvsubsection{Virtualizzazione e orchestrazione}
\cvskill{Docker (Swarm, Compose), Podman, Buildah, Kubernetes, Skaffold, LXC, VirtualBox, systemd-nspawn}

\vspace{-2mm}
\cvsubsection{Automazione e documentazione}
\cvskill{GitLab CI/CD, Make, Maven, Gradle, CMake}

\vspace{-2mm}
\cvsubsection{Database e simili}
\cvskill{MariaDB, SQLite, MySQL, PostgreSQL, Redis, ArangoDB, MongoDB}

\vspace{-2mm}
\cvsubsection{Protocolli o standard}
\cvskill{HTTP, POSIX, OpenAPI, TCP, UDP, SSH, PGP/GPG, ANSI C, UML, JSON, YAML, XML Unicode}

\vspace{-2mm}
\cvsubsection{Miscellanea}
\cvskill{REST, Assembly MIPS, Systemd, GNU Groff, \LaTeX, AsciiDoc, Markdown, PlantUML}

% -----------------------
% Certifications section.
% -----------------------
\cvsection{Certificazioni}

\cvaward{Attestato di partecipazione CyberChallenge.IT 2021}{CyberChallenge.IT}{giu 2021}

\cvaward{Java SE 8 Programmer}{Oracle}{lug 2018}

% ---------------------
% Volunteering section.
% ---------------------
\cvsection{Volontariato}

\cvinvolve{Software libero e open source (FOSS)}{2015 -- IN CORSO}{}
\begin{itemize}\setlength\itemsep{0.2em}
\item Contributore a \textbf{Fedora Project} per la localizzazione in italiano tramite l'uso della piattaforma Weblate.
\item Contributore a \textbf{Xfce Desktop Environment} nello sviluppo, testing e rilascio del software scritto in linguaggio \textbf{C} e \textbf{Python} con utilizzo di librerie GLib/GTK.
\item Contributore a una moltitudine di progetti minori, sia personali che della comunità, in tutte le dinamiche: documentazione, sviluppo, testing, supporto, localizzazione, licenze software, lavoro in gruppo e singolo e altro.
\end{itemize}

\cvinvolve{PCOfficina}{2022 -- IN CORSO}{}
\begin{itemize}\setlength\itemsep{0.2em}
\item Volontario nell'associazione con l'obiettivo di riduzione dei rifiuti tecnologici informatici, diffusione e divulgazione conoscenze informatiche e abbassamento del ``digital divide'' tramite il recupero hardware e software di computer con successiva donazione.
\end{itemize}



\end{paracol}

\end{document}
